% Zinsrechnung – bank/zinsrechnung/ (zusammenbau v0.4, Heft)
\einheitenkopf[t4]{Zinsrechnung}
\setcounter{aufgabe}{34}

% Anlauf (ohne Prüfkennung)
\begin{aufgabe}{Zinsen berechnen – Anlauf}
\begin{teile}
\teil Lena legt 850\,€ für ein Jahr auf ein Sparbuch, die Bank zahlt 2\,\% Zinsen. Wie viel Euro Zinsen bekommt sie? Was ist gesucht? Kreuze an, rechne nicht. \\ \kreuz{das Kapital (Grundwert)} \\ \kreuz{der Zinssatz (Prozentsatz)} \\ \kreuz{die Zinsen (Prozentwert)}
\teil Sparbuch: 300\,€ zu 2\,\% – wie viel Euro Zinsen nach einem Jahr? \leerfeld[€]
\teil Rechne über 1\,\%: 600\,€ zu 4\,\% – wie viel Euro Zinsen im Jahr? \leerfeld[€]
\end{teile}
\end{aufgabe}

\begin{aufgabe}{Zinsen berechnen – Prüfungsaufgaben}
\begin{teile}
\teil Ein Sparbetrag von 2\,600\,€ wird ein Jahr lang mit 2\,\% pro Jahr verzinst. Wie viel Euro Zinsen gibt es nach einem Jahr? \hfill (P10 2015 OS) \\ \leerfeld[€]
\teil Eine Geldanlage von 4\,500\,€ wird ein Jahr lang mit 3\,\% pro Jahr verzinst. Wie viel Euro Zinsen gibt es nach einem Jahr? \hfill (P10 2015 OS) \\ \leerfeld[€]
\teil Ein Sparguthaben von 15\,000\,€ bringt in einem Jahr 270\,€ Zinsen. Gib den Zinssatz an. \hfill (P10 2014 OS) \\ \leerfeld[\%]
\teil Ein Sparguthaben von 2\,500\,€ bringt in einem Jahr 65\,€ Zinsen. Gib den Zinssatz an. \hfill (P10 2014 OS) \\ \leerfeld[\%]
\teil Auf ein Sparbuch werden jedes Jahr 250\,€ eingezahlt. Weise nach, dass die Bank 3\,\% Zinsen pro Jahr zahlt. \hfill (P10 2014 OS)

\sachtabelle{lrrr}{Jahr & Zinsen & Einzahlung & Guthaben}{1 & – & 250,00\,€ & 250,00\,€\\ 2 & 7,50\,€ & 250,00\,€ & 507,50\,€}
\teil Auf ein Sparbuch werden jedes Jahr 350\,€ eingezahlt. Weise nach, dass die Bank 2\,\% Zinsen pro Jahr zahlt. \hfill (P10 2014 OS)

\sachtabelle{lrrr}{Jahr & Zinsen & Einzahlung & Guthaben}{1 & – & 350,00\,€ & 350,00\,€\\ 2 & 7,00\,€ & 350,00\,€ & 707,00\,€}
\end{teile}
\end{aufgabe}

% Anlauf (ohne Prüfkennung)
\begin{aufgabe}{Zinseszins – Anlauf}
\begin{teile}
\teil Auf einem Konto liegen am Anfang 600\,€, nach einem Jahr mit Zinsen 612\,€. Von welchem Betrag werden die Zinsen im zweiten Jahr berechnet? Kreuze an, rechne nicht. \\ \kreuz{vom neuen Guthaben 612\,€} \\ \kreuz{vom Startguthaben 600\,€}
\teil 900\,€ zu 3\,\%, die Zinsen werden am Jahresende gutgeschrieben – neues Guthaben? \leerfeld[€]
\teil 2\,000\,€ zu 3\,\%, die Zinsen bleiben auf dem Konto. Rechne Jahr für Jahr – Guthaben nach zwei Jahren? \leerfeld[€]
\end{teile}
\end{aufgabe}

\begin{aufgabe}{Zinseszins – Prüfungsaufgaben}
\begin{teile}
\teil Für den Führerschein zahlt Lukas jedes Jahr 350\,€ auf ein Konto ein, der Zinssatz ist 3\,\%. Vervollständige die Tabelle und kontrolliere mit Jahr 4. \hfill (P10 2014 OS) \\

\sachtabelle{lrrr}{Jahr & Zinsen & Einzahlung & Guthaben}{1 & – & 350,00\,€ & 350,00\,€\\ 2 & 10,50\,€ & 350,00\,€ & 710,50\,€\\ 3 & 21,32\,€ & 350,00\,€ & \leerzelle\\ 4 & \leerzelle & 350,00\,€ & 1\,464,27\,€}
Guthaben \leerfeld[€] , Zinsen \leerfeld[€]
\teil Für Emma werden jedes Jahr 180\,€ auf ein Sparkonto eingezahlt, der Zinssatz ist 2\,\%. Vervollständige die Tabelle und kontrolliere mit Jahr 4. \hfill (P10 2014 OS) \\

\sachtabelle{lrrr}{Jahr & Zinsen & Einzahlung & Guthaben}{1 & – & 180,00\,€ & 180,00\,€\\ 2 & 3,60\,€ & 180,00\,€ & 363,60\,€\\ 3 & 7,27\,€ & 180,00\,€ & \leerzelle\\ 4 & \leerzelle & 180,00\,€ & 741,89\,€}
Guthaben \leerfeld[€] , Zinsen \leerfeld[€]
\teil Auf einen Bausparvertrag werden 4\,000\,€ angelegt, 2\,\% pro Jahr. Die Zinsen bleiben stehen und werden mitverzinst. Berechne das Guthaben nach 6 Jahren. \hfill (P10 2014 OS) \\ \leerfeld[€]
\teil In einen Sparbrief werden 2\,500\,€ angelegt, 3\,\% pro Jahr. Die Zinsen bleiben stehen und werden mitverzinst. Berechne das Guthaben nach 7 Jahren. \hfill (P10 2014 OS) \\ \leerfeld[€]
\end{teile}
\end{aufgabe}
